\begin{abstract}
Demand-responsive transit (DRT) offers flexible mobility in low-density urban areas,
but door-to-door service is operationally costly and difficult to scale.
Meeting-point-based DRT reduces vehicle detours by consolidating pickups and dropoffs
at pre-defined locations, yet existing work assigns meeting points only on the pickup
side and ignores passenger heterogeneity in mode choice.
We address this gap by proposing a many-to-many DRT model with bidirectional meeting
point sets ($M_r^P$ for pickup, $M_r^D$ for dropoff) and passenger choice modeled
via single-offer binary logit acceptance with three passenger types and an outside option.
The model is formulated as a coupled three-layer system — service offer generation,
passenger response, and dynamic dispatch — and solved with an ex-post MILP routing
diagnostic for small instances and a rolling horizon Adaptive
Large Neighborhood Search (ALNS) heuristic for large-scale deployment.
Numerical experiments on synthetic scenarios calibrated to Chinese suburban conditions
show that, at equal served share (endogenous matched-coverage comparison), the full bidirectional model
achieves significantly lower per-trip vehicle utilization than the
DoorToDoor\,(capped) baseline.
Under the unconstrained comparison (FullModel 18.3\%
vs.\ DoorToDoor 61.0\% served share), FullModel achieves 15.1\,vkm/trip versus
21.3\,vkm/trip (29.1\% improvement). For reference, a post-hoc matched-coverage
estimate confirms the efficiency advantage is not an artifact of lower coverage.
Equity analysis across passenger types reveals a Gini coefficient of 0.1216 in
acceptance rates: time-sensitive passengers are systematically disadvantaged
(14.4\% acceptance) relative to price-sensitive (25.4\%) and walk-sensitive (20.2\%)
passengers.
Based on these findings, we derive five scenario-specific policy implications for
DRT deployment in low-density Chinese urban areas.
\end{abstract}

\begin{keyword}
demand-responsive transit \sep meeting points \sep passenger choice \sep ALNS \sep rolling horizon \sep equity
\end{keyword}
